\chapter{2004年湖南卷（理）}
\section{单选题}
\begin{problem}
复数 \(\left(1 + \frac{1}{\i}\right)^4\) 的值是（\quad）

\choices
    {\(4\i\)}
    {\(-4\i\)}
    {\(4\)}
    {\(-4\)}
\end{problem}

\begin{answer}
D
\end{answer}

\begin{problem}
如果双曲线 \(\frac{x^2}{13} - \frac{y^2}{12} = 1\) 上一点 \(P\) 到右焦点的距离等于 \(\sqrt{13}\)，那么点 \(P\) 到右准线的距离是（\quad）

\choices
    {\(\frac{13}{5}\)}
    {\(13\)}
    {\(5\)}
    {\(\frac{5}{13}\)}
\end{problem}

\begin{answer}
A
\end{answer}

\begin{problem}
设 \(f^{-1}(x)\) 是函数 \(f(x)=\log_{2}(x+1)\) 的反函数，若 \(\left[1+f^{-1}(a)\right]\left[1+f^{-1}(b)\right]=8\)，则 \(f(a-b)\) 的值为（\quad）

\choices
    {\(1\)}
    {\(2\)}
    {\(3\)}
    {\(\log_{2}3\)}
\end{problem}

\begin{answer}
题干条件不足，无法确定.
\end{answer}

\begin{solution}
由 \(f(x)=\log_{2}(x+1)\) 可得 \(f^{-1}(x)=2^x-1\)，所以
\[
\left[1+f^{-1}(a)\right]\left[1+f^{-1}(b)\right]=2^a2^b=2^{a+b}=8=2^3,
\]
从而 \(a+b=3\). 但是，这个条件不能唯一确定 \(a-b\)，因而不能唯一确定 \(f(a-b)\). 例如，取 \((a,b)=(2,1)\)，则 \(f(a-b)=f(1)=1\)；取 \((a,b)=(3,0)\)，则 \(f(a-b)=f(3)=2\). 这两组取值都满足题设条件，但所得函数值不同；取 \((a,b)=(1,2)\) 时还有 \(a-b=-1\)，此时 \(f(a-b)\) 无定义. 因此按原题题干，\(f(a-b)\) 的值不能确定，四个选项中没有唯一正确的选项.
\end{solution}

\begin{problem}
把正方形 \(ABCD\) 沿对角线 \(AC\) 折起，当以 \(A\)，\(B\)，\(C\)，\(D\) 四点为顶点的三棱锥体积最大时，直线 \(BD\) 和平面 \(ABC\) 所成的角的大小为（\quad）

\choices
    {\(90\)}
    {\(60\)}
    {\(45\)}
    {\(30\)}
\end{problem}

\begin{answer}
C
\end{answer}

\begin{problem}
某公司在甲，乙，丙，丁四个地区分别有 \(150\text{ 个}\)，\(120\text{ 个}\)，\(180\text{ 个}\)，\(150\text{ 个}\)销售点. 公司为了调查产品销售的情况，需从这 \(600\text{ 个}\)销售点中抽取一个容量为 \(100\) 的样本，记这项调查为\(\circled{1}\)；在丙地区中有 \(20\text{ 个}\)特大型销售点，要从中抽取 \(7\text{ 个}\)调查其销售收入和售后服务情况，记这项调查为\(\circled{2}\). 则完成\(\circled{1}\)，\(\circled{2}\)这两项调查宜采用的抽样方法依次是（\quad）

\choices
    {分层抽样，系统抽样法}
    {分层抽样，简单随机抽样法}
    {系统抽样法，分层抽样法}
    {简单随机抽样法，分层抽样法}
\end{problem}

\begin{answer}
B
\end{answer}

\begin{problem}
设函数 \(f(x)=\begin{cases}
x^2+bx+c, & x\leqslant 0,\\
2, & x>0.
\end{cases}\) 若 \(f(-4)=f(0)\)，\(f(-2)=-2\)，则关于 \(x\) 的方程 \(f(x)=x\) 的解的个数为（\quad）

\choices
    {\(1\)}
    {\(2\)}
    {\(3\)}
    {\(4\)}
\end{problem}

\begin{answer}
C
\end{answer}

\begin{problem}
设 \(a>0\)，\(b>0\)，则以下不等式中不恒成立的是（\quad）

\choices
    {\(\left(a+b\right)\left(\frac{1}{a}+\frac{1}{b}\right)\geqslant 4\)}
    {\(a^3+b^3\geqslant 2ab^2\)}
    {\(a^2+b^2+2\geqslant 2a+2b\)}
    {\(\sqrt{|a-b|}\geqslant \sqrt{a}-\sqrt{b}\)}
\end{problem}

\begin{answer}
B
\end{answer}

\begin{problem}
数列 \(\{a_n\}\) 中，\(a_1=\frac{1}{5}\)，\(a_n+a_{n+1}=\frac{6}{5^{n+1}}\)，\(n\in \N^{*}\)，则 \(\lim\limits_{n\to\infty}(a_1+a_2+\cdots+a_n)=\)（\quad）

\choices
    {\(\frac{2}{5}\)}
    {\(\frac{2}{7}\)}
    {\(\frac{1}{4}\)}
    {\(\frac{4}{25}\)}
\end{problem}

\begin{answer}
C
\end{answer}

\begin{problem}
设集合 \(U=\{(x,y)\mid x\in\R, y\in\R\}\)，\(A=\{(x,y)\mid 2x-y+m>0\}\)，\(B=\{(x,y)\mid x+y-n\leqslant 0\}\)，那么点 \(P(2,3)\in A\cap(\complement_{U}B)\) 的充要条件是（\quad）

\choices
    {\(m>-1, n<5\)}
    {\(m<-1\)，\(n<5\)}
    {\(m>-1\)，\(n>5\)}
    {\(m<-1\)，\(n>5\)}
\end{problem}

\begin{answer}
A
\end{answer}

\begin{problem}
从正方体的八个顶点中任取三个点为顶点作三角形，其中直角三角形的个数为（\quad）

\choices
    {\(56\)}
    {\(52\)}
    {\(48\)}
    {\(40\)}
\end{problem}

\begin{answer}
C
\end{answer}

\begin{problem}
农民收入由工资性收入和其它收入两部分构成. 2003年某地区农民人均收入为 \(3150\text{ 元}\)（其中工资性收入为 \(1800\text{ 元}\)，其它收入为 \(1350\text{ 元}\)），预计该地区自 2004 年起的 \(5\text{ 年}\)内，农民的工资性收入将以每年 \(6\%\) 的年增长率增长，其它收入每年增加 \(160\text{ 元}\). 根据以上数据，2008 年该地区农民人均收入介于（\quad）

\choices
    {\(4200\text{ 元}\sim 4400\text{ 元}\)}
    {\(4400\text{ 元}\sim 4600\text{ 元}\)}
    {\(4600\text{ 元}\sim 4800\text{ 元}\)}
    {\(4800\text{ 元}\sim 5000\text{ 元}\)}
\end{problem}

\begin{answer}
B
\end{answer}

\begin{problem}
设 \(f(x)\)，\(g(x)\) 分别是定义在 \(\R\) 上的奇函数和偶函数，当 \(x<0\) 时，\(f'(x)g(x)+f(x)g'(x)>0\). 且 \(g(-3)=0\)，则不等式 \(f(x)g(x)<0\) 的解集是（\quad）

\choices
    {\((-3,0)\cup(3,+\infty)\)}
    {\((-3,0)\cup(0,3)\)}
    {\((-\infty,-3)\cup(3,+\infty)\)}
    {\((-\infty,-3)\cup(0,3)\)}
\end{problem}

\begin{answer}
D
\end{answer}

\section{填空题}
\begin{problem}
已知向量 \(\bs{a}=(\cos\theta,\sin\theta)\)，向量 \(\bs{b}=(\sqrt{3},-1)\)，则 \(\left|2\bs{a}-\bs{b}\right|\) 的最大值是\fillinblank{}.
\end{problem}

\begin{answer}
\(4\)
\end{answer}

\begin{problem}
同时抛掷两枚相同的均匀硬币，随机变量 \(\xi=1\) 表示结果中有正面向上，\(\xi=0\) 表示结果中没有正面向上，则 \(E(\xi)=\)\fillinblank{}.
\end{problem}

\begin{answer}
\(0.75\)
\end{answer}

\begin{problem}
若 \(\left(x^{3}+\frac{1}{x\sqrt{x}}\right)^n\) 展开式中的常数项为 84，则 \(n=\)\fillinblank{}.
\end{problem}

\begin{answer}
\(9\)
\end{answer}

\begin{problem}
设 \(F\) 是椭圆 \(\frac{x^2}{7}+\frac{y^2}{6}=1\) 的右焦点，且椭圆上至少有 \(21\) 个不同的点 \(P_i\)（\(i=1,2,3,\dots\)），使 \(\left|FP_1\right|\)，\(\left|FP_2\right|\)，\(\left|FP_3\right|\)，\ldots 组成公差为 \(d\) 的等差数列，则 \(d\) 的取值范围为\fillinblank{}.
\end{problem}

\begin{answer}
\(\left[-\frac{1}{10},0\right)\cup\left(0,\frac{1}{10}\right]\)
\end{answer}

\section{解答题}
\begin{problem}
已知 \(\sin\left(\frac{\pi}{4}+2\alpha\right)\cdot\sin\left(\frac{\pi}{4}-2\alpha\right)=\frac{1}{4}\)，（\(\alpha\in\left(\frac{\pi}{4},\frac{\pi}{2}\right)\)），求 \(2\sin^2\alpha+\tan\alpha-\cot\alpha-1\) 的值.
\end{problem}

\begin{solution}
由积化和差公式，原条件等价于
\[
\sin^2\frac{\pi}{4}-\sin^2 2\alpha=\frac{1}{4},
\]
所以 \(\sin^2 2\alpha=\frac{1}{4}\). 因为 \(2\alpha\in\left(\frac{\pi}{2},\pi\right)\)，故 \(\sin 2\alpha>0\)，从而 \(\sin 2\alpha=\frac{1}{2}\)，于是 \(2\alpha=\frac{5\pi}{6}\)，即 \(\alpha=\frac{5\pi}{12}\). 又由 \(2\sin^2\alpha-1=-\cos 2\alpha\)，有
\[
\tan\alpha-\cot\alpha=\frac{\sin^2\alpha-\cos^2\alpha}{\sin\alpha\cos\alpha}=-\frac{2\cos 2\alpha}{\sin 2\alpha}.
\]
代入 \(\sin 2\alpha=\frac{1}{2}\)，\(\cos 2\alpha=-\frac{\sqrt{3}}{2}\)，得
\[
2\sin^2\alpha+\tan\alpha-\cot\alpha-1=\frac{\sqrt{3}}{2}+2\sqrt{3}=\frac{5\sqrt{3}}{2}.
\]
\end{solution}

\begin{problem}
如图，在底面是菱形的四棱锥 \(P-ABCD\) 中，\(\angle ABC=60^{\circ}\)，\(PA=AC=a\)，\(PB=PD=\sqrt{2}a\)，点 \(E\) 在 \(PD\) 上，且 \(PE:ED=2:1\).

\begin{enumerate}
\item 证明：\(PA\perp\) 平面 \(ABCD\)；
\item 求以 \(AC\) 为棱，\(EAC\) 与 \(DAC\) 为面的二面角 \(\theta\) 的大小；
\item 在棱 \(PC\) 上是否存在一点 \(F\)，使 \(BF\parallel\) 平面 \(AEC\)？ 证明你的结论.
\end{enumerate}

\bitmapfigure[max width=0.45\linewidth]{img/2004/hunan_science/q18_fig1.png}
\end{problem}

\begin{solution}
\begin{enumerate}
\item 因为 \(ABCD\) 是菱形且 \(\angle ABC=60^\circ\)，所以三角形 \(ABC\) 和三角形 \(ADC\) 都是边长为 \(a\) 的等边三角形，即 \(AB=AD=AC=a\). 在三角形 \(PAB\) 中，\(PA^2+AB^2=a^2+a^2=2a^2=PB^2\)，所以 \(PA\perp AB\). 同理，在三角形 \(PAD\) 中，\(PA^2+AD^2=PD^2\)，所以 \(PA\perp AD\). 由于 \(AB\) 与 \(AD\) 是平面 \(ABCD\) 内的相交直线，故 \(PA\perp\) 平面 \(ABCD\).
\item 过 \(E\) 作 \(EG\parallel PA\)，交 \(AD\) 于 \(G\). 因为 \(PE:ED=2:1\)，所以 \(DE:DP=1:3\)，由相似三角形得 \(DG:DA=1:3\)，\(EG:PA=1:3\)，从而 \(DG=\frac{a}{3}\)，\(EG=\frac{a}{3}\). 由于三角形 \(ADC\) 为等边三角形，且 \(G\in AD\)，作 \(GH\perp AC\) 于 \(H\)，则 \(AG=\frac{2a}{3}\)，\(GH=AG\sin60^\circ=\frac{\sqrt{3}a}{3}\). 又 \(EG\perp\) 平面 \(ABCD\)，所以 \(EG\perp GH\)，且 \(EH\perp AC\). 因此 \(\angle EHG\) 是二面角 \(EAC-DAC\) 的平面角，并且
\[
\tan\theta=\frac{EG}{GH}=\frac{a/3}{\sqrt{3}a/3}=\frac{1}{\sqrt{3}}.
\]
所以 \(\theta=30^\circ\).
\item 建立空间直角坐标系：取 \(A\) 为原点，令 \(AD\) 为 \(y\) 轴，\(AP\) 为 \(z\) 轴，取 \(x\) 轴使其垂直于平面 \(PAD\). 可设 \(A(0,0,0)\)，\(D(0,a,0)\)，\(P(0,0,a)\)，\(B\left(\frac{\sqrt{3}a}{2},-\frac{a}{2},0\right)\)，\(C\left(\frac{\sqrt{3}a}{2},\frac{a}{2},0\right)\). 由 \(PE:ED=2:1\)，得 \(E\left(0,\frac{2a}{3},\frac{a}{3}\right)\). 取 \(F\) 为棱 \(PC\) 的中点，则 \(F\left(\frac{\sqrt{3}a}{4},\frac{a}{4},\frac{a}{2}\right)\). 因此
\[
\overrightarrow{BF}=\left(-\frac{\sqrt{3}a}{4},\frac{3a}{4},\frac{a}{2}\right).
\]
并且 \(\overrightarrow{AE}=\left(0,\frac{2a}{3},\frac{a}{3}\right)\)，\(\overrightarrow{AC}=\left(\frac{\sqrt{3}a}{2},\frac{a}{2},0\right)\).
\[
\overrightarrow{BF}=\frac{3}{2}\overrightarrow{AE}-\frac{1}{2}\overrightarrow{AC}.
\]
所以 \(BF\) 的方向向量是平面 \(AEC\) 内两个相交方向 \(AE\)、\(AC\) 的线性组合. 又若 \(B\) 在平面 \(AEC\) 内，则存在实数 \(u,v\)，使 \(\overrightarrow{AB}=u\overrightarrow{AE}+v\overrightarrow{AC}\). 比较上式各坐标可得 \(v=1\)、\(u=0\)，但此时右端的 \(y\) 坐标为 \(\frac{a}{2}\)，而 \(B\) 的 \(y\) 坐标为 \(-\frac{a}{2}\)，矛盾. 所以 \(B\) 不在平面 \(AEC\) 内，故直线 \(BF\) 平行于平面 \(AEC\). 因而存在这样的点 \(F\)，且 \(F\) 是 \(PC\) 的中点.
\end{enumerate}
\end{solution}

\begin{problem}
甲，乙，丙三台机床各自独立地加工同一种零件，已知甲机床加工的零件是一等品而乙机床加工的零件不是一等品的概率为 \(\frac{1}{4}\). 乙机床加工的零件是一等品而丙机床加工的零件不是一等品的概率为 \(\frac{1}{12}\). 甲，丙两台机床加工的零件都是一等品的概率为 \(\frac{2}{9}\).

\begin{enumerate}
\item 分别求甲，乙，丙三台机床各自加工零件是一等品的概率；
\item 从甲，乙，丙加工的零件中各取一个检验，求至少有一个一等品的概率.
\end{enumerate}
\end{problem}

\begin{solution}
设甲、乙、丙机床加工的零件为一等品的概率分别为 \(p\)、\(q\)、\(r\). 由各机床独立加工以及题设条件，得
\[
\begin{cases}
p(1-q)=\dfrac{1}{4},\\
q(1-r)=\dfrac{1}{12},\\
pr=\dfrac{2}{9}.
\end{cases}
\]
由第一式得 \(q=1-\frac{1}{4p}=\frac{4p-1}{4p}\)，再由第二式得
\[
r=1-\frac{1}{12q}=\frac{11p-3}{12p-3}.
\]
代入 \(pr=\frac{2}{9}\)，得
\[
\frac{p(11p-3)}{12p-3}=\frac{2}{9},
\qquad
33p^2-17p+2=0.
\]
所以 \(p=\frac{1}{3}\) 或 \(p=\frac{2}{11}\). 由 \(q=1-\frac{1}{4p}>0\) 得 \(p>\frac{1}{4}\)，舍去 \(p=\frac{2}{11}\)，从而 \(p=\frac{1}{3}\)，\(q=\frac{1}{4}\)，\(r=\frac{2}{3}\).
至少有一个一等品的概率为
\[
1-(1-p)(1-q)(1-r)=1-\frac{2}{3}\cdot\frac{3}{4}\cdot\frac{1}{3}=\frac{5}{6}.
\]
\end{solution}

\begin{problem}
已知函数 \(f(x)=x^{2}\e^{ax}\)，其中 \(a\leqslant 0\)，\(\e\) 为自然对数的底数.

\begin{enumerate}
\item 讨论函数 \(f(x)\) 的单调性；
\item 求函数 \(f(x)\) 在区间 \([0,1]\) 上的最大值.
\end{enumerate}
\end{problem}

\begin{solution}
\begin{enumerate}
\item
\[
f'(x)=2x\e^{ax}+ax^2\e^{ax}=x\e^{ax}(2+ax).
\]
因为 \(\e^{ax}>0\)，先看 \(a=0\). 此时 \(f'(x)=2x\)，所以 \(f(x)\) 在 \((-\infty,0)\) 上单调递减，在 \((0,+\infty)\) 上单调递增.

当 \(a<0\) 时，临界点为 \(x=0\) 和 \(x=-\frac{2}{a}\). 当 \(x<0\) 时，\(x<0\) 且 \(2+ax>0\)，所以 \(f'(x)<0\)；当 \(0<x<-\frac{2}{a}\) 时，\(f'(x)>0\)；当 \(x>-\frac{2}{a}\) 时，\(f'(x)<0\). 因此 \(f(x)\) 在 \((-\infty,0)\) 上单调递减，在 \(\left(0,-\frac{2}{a}\right)\) 上单调递增，在 \(\left(-\frac{2}{a},+\infty\right)\) 上单调递减.
\item 当 \(-2\leqslant a\leqslant0\) 时，\(-\frac{2}{a}\geqslant1\)（\(a=0\) 时直接由单调性得结论），故 \(f(x)\) 在 \([0,1]\) 上单调递增，最大值为 \(f(1)=\e^a\). 当 \(a<-2\) 时，\(0<-\frac{2}{a}<1\)，最大值在 \(x=-\frac{2}{a}\) 处取得，且
\[
f\left(-\frac{2}{a}\right)=\frac{4}{a^2}\e^{-2}.
\]
所以区间 \([0,1]\) 上的最大值为
\[
\begin{cases}
\e^a,&-2\leqslant a\leqslant0,\\
\dfrac{4}{a^2\e^2},&a<-2.
\end{cases}
\]
\end{enumerate}
\end{solution}

\begin{problem}
如图，过抛物线 \(x^{2}=4y\) 的对称轴上任一点 \(P(0,m)\)（\(m>0\)）作直线与抛物线交于 \(A\)，\(B\) 两点，点 \(Q\) 是点 \(P\) 关于原点的对称点.

\begin{enumerate}
\item 设点 \(P\) 分有向线段 \(\overrightarrow{AB}\) 所成的比为 \(\lambda\)，证明：\(\overrightarrow{QP}\perp(\overrightarrow{QA}-\lambda\overrightarrow{QB})\)；
\item 设直线 \(AB\) 的方程是 \(x-2y+12=0\)，过 \(A\)，\(B\) 两点的圆 \(C\) 与抛物线在点 \(A\) 处有共同的切线，求圆 \(C\) 的方程.
\end{enumerate}

\bitmapfigure[max width=0.34\linewidth]{img/2004/hunan_science/q21_fig1.png}
\end{problem}

\begin{solution}
\begin{enumerate}
\item 设直线 \(AB\) 的方程为 \(y=tx+m\)，设 \(A(x_1,y_1)\)、\(B(x_2,y_2)\). 联立 \(y=\frac{x^2}{4}\) 得
\[
x^2-4tx-4m=0,
\]
所以 \(x_1+x_2=4t\)，\(x_1x_2=-4m\)，且 \(y_i=\frac{x_i^2}{4}\)（\(i=1,2\)）.
由有向线段比的定义，
\[
\lambda=\frac{\overline{AP}}{\overline{PB}}=-\frac{x_1}{x_2}.
\]
又 \(Q=(0,-m)\)，所以 \(\overrightarrow{QP}=(0,2m)\). 由 \(m=-\frac{x_1x_2}{4}\)，有
\[
y_1+m=\frac{x_1(x_1-x_2)}{4},\qquad y_2+m=-\frac{x_2(x_1-x_2)}{4}.
\]
因此 \(\overrightarrow{QA}-\lambda\overrightarrow{QB}\) 的纵坐标为
\[
(y_1+m)-\lambda(y_2+m)=\frac{x_1(x_1-x_2)}{4}-\left(-\frac{x_1}{x_2}\right)\left(-\frac{x_2(x_1-x_2)}{4}\right)=0.
\]
故 \(\overrightarrow{QA}-\lambda\overrightarrow{QB}\) 是水平向量，从而
\[
\overrightarrow{QP}\cdot\left(\overrightarrow{QA}-\lambda\overrightarrow{QB}\right)=0.
\]
所以 \(\overrightarrow{QP}\perp\left(\overrightarrow{QA}-\lambda\overrightarrow{QB}\right)\).
\item 由直线 \(AB\) 的方程 \(x-2y+12=0\)，得 \(y=\frac{x}{2}+6\). 联立抛物线方程，得
\[
\frac{x^2}{4}=\frac{x}{2}+6,
\qquad
x^2-2x-24=0,
\]
所以 \(A(6,9)\)、\(B(-4,4)\). 抛物线写成 \(y=\frac{x^2}{4}\)，其在 \(A\) 处切线斜率为 \(\frac{x}{2}\big|_{x=6}=3\)，故切线方程为 \(y-9=3(x-6)\)，即 \(3x-y-9=0\).

设圆 \(C\) 的方程为 \(x^2+y^2+ux+vy+w=0\). 因为 \(A\)、\(B\) 在圆上，且圆在 \(A\) 处的切线与直线 \(3x-y-9=0\) 重合，故
\[
\begin{cases}
117+6u+9v+w=0,\\
32-4u+4v+w=0,\\
u+3v+66=0.
\end{cases}
\]
第三式来自圆在 \(A\) 处的法向量 \((12+u,18+v)\) 与 \((3,-1)\) 平行. 解得 \(u=3\)、\(v=-23\)、\(w=72\). 所以圆 \(C\) 的方程为
\[
x^2+y^2+3x-23y+72=0.
\]
\end{enumerate}
\end{solution}

\begin{problem}
如图，直线 \(l_1:y=kx+1-k\)（\(k\neq 0\)，\(k\neq\pm\frac{1}{2}\)）与 \(l_2:y=\frac{1}{2}x+\frac{1}{2}\) 相交于点 \(P\). 直线 \(l_1\) 与 \(x\) 轴交于点 \(P_1\)，过点 \(P_1\) 作 \(x\) 轴的垂线交直线 \(l_2\) 于点 \(Q_1\)，过点 \(Q_1\) 作 \(y\) 轴的垂线交直线 \(l_1\) 于点 \(P_2\)，过点 \(P_2\) 作 \(x\) 轴的垂线交直线 \(l_2\) 于点 \(Q_2\)，\(\cdots\)，这样一直作下去，可得到一系列点 \(P_1\)，\(Q_1\)，\(P_2\)，\(Q_2\)，\(\cdots\). 点 \(P_n\)（\(n=1,2,\cdots\)）的横坐标构成数列 \(\{x_n\}\).

\begin{enumerate}
\item 证明：\(x_{n+1}-1=\frac{1}{2k}(x_n-1)\)，\(n\in\N^{*}\)；
\item 求数列 \(\{x_n\}\) 的通项公式；
\item 比较 \(2\left|PP_n\right|^{2}\) 与 \(4k^{2}\left|PP_1\right|^{2}+5\) 的大小.
\end{enumerate}

\bitmapfigure[max width=0.42\linewidth]{img/2004/hunan_science/q22_fig1.png}
\end{problem}

\begin{solution}
\begin{enumerate}
\item 由 \(P_1\) 在 \(x\) 轴上，得 \(P_1=(x_1,0)\). 过 \(P_n=(x_n,kx_n+1-k)\) 作 \(x\) 轴的垂线，与 \(l_2\) 的交点为
\[
Q_n=\left(x_n,\frac{x_n}{2}+\frac{1}{2}\right).
\]
过 \(Q_n\) 作 \(y\) 轴的垂线，与 \(l_1\) 交于 \(P_{n+1}\)，故
\[
kx_{n+1}+1-k=\frac{x_n}{2}+\frac{1}{2}.
\]
整理得
\[
x_{n+1}=1+\frac{x_n-1}{2k},
\]
即 \(x_{n+1}-1=\frac{1}{2k}(x_n-1)\).
\item 由 \(P_1\) 在 \(x\) 轴上，\(kx_1+1-k=0\)，所以 \(x_1=1-\frac{1}{k}\). 由上一问，数列 \(\{x_n-1\}\) 是首项 \(-\frac{1}{k}\)、公比 \(\frac{1}{2k}\) 的等比数列，从而
\[
x_n-1=-\frac{1}{k}\left(\frac{1}{2k}\right)^{n-1},
\qquad
x_n=1-\frac{1}{k(2k)^{n-1}}.
\]
\item 由 \(P=(1,1)\)，且 \(P_n\in l_1\)，有
\[
y_n-1=k(x_n-1)=-\left(\frac{1}{2k}\right)^{n-1}.
\]
因此
\[
\left|PP_n\right|^2=(x_n-1)^2+(y_n-1)^2=\frac{1+k^2}{k^2(2k)^{2n-2}},
\]
而 \(\left|PP_1\right|^2=\frac{1+k^2}{k^2}\). 于是
\[
4k^2\left|PP_1\right|^2+5=4k^2+9,
\]
且当 \(n=1\) 时两待比较量之差为
\[
2\left|PP_1\right|^2-\left(4k^2\left|PP_1\right|^2+5\right)=\frac{(1-4k^2)(2+k^2)}{k^2}.
\]
若 \(\left|k\right|<\frac{1}{2}\)，则 \(0<4k^2<1\)，所以 \(2\left|PP_n\right|^2\geq 2\left|PP_1\right|^2>4k^2\left|PP_1\right|^2+5\). 若 \(\left|k\right|>\frac{1}{2}\)，则 \(4k^2>1\)，所以 \(2\left|PP_n\right|^2\leq 2\left|PP_1\right|^2<4k^2\left|PP_1\right|^2+5\). 因而
\[
\begin{cases}
2\left|PP_n\right|^2>4k^2\left|PP_1\right|^2+5,&\left|k\right|<\frac{1}{2},\\
2\left|PP_n\right|^2<4k^2\left|PP_1\right|^2+5,&\left|k\right|>\frac{1}{2}.
\end{cases}
\]
\end{enumerate}
\end{solution}
